% BANKS-SOURCE: pdf=99 print=89 kind=implicit-exercise id=I-CH07-005
\begin{exercise}[The composite $H$-connection]{I-CH07-005}
Let $H\subset G$ and choose an $\operatorname{Ad}(H)$-invariant decomposition
$\mathfrak g=\mathfrak h\oplus\mathfrak m$.  Let $t_m$ be a basis of
$\mathfrak h$ orthonormal for a chosen invariant bilinear form, denoted by
$\Tr$, so that $\Tr(t_mt_n)=\delta_{mn}$.  For anti-Hermitian matrices this
notation may include the conventional minus sign relative to the ordinary
matrix trace.  A coset field is represented by $g(x)\in G$ with the local
right equivalence $g(x)\sim g(x)h(x)$, $h(x)\in H$.  Define
\[
 A_\mu=P_{\mathfrak h}(g^{-1}\partial_\mu g)
       =t_m A_\mu^m,
 \qquad
 A_\mu^m=\Tr(t_m g^{-1}\partial_\mu g).
\]
\begin{enumerate}[label=(\alph*)]
\item Show that under $g\mapsto gh$,
\[
 A_\mu\mapsto h^{-1}(A_\mu+\partial_\mu)h
 =h^{-1}A_\mu h+h^{-1}\partial_\mu h,
\]
so that $A_\mu$ is an $H$-connection.

\item In an anti-Hermitian generator convention, give the infinitesimal
component transformation and verify directly that
$D_\mu g=\partial_\mu g-gA_\mu$ transforms as
$D_\mu g\mapsto(D_\mu g)h$.  Check that the curvature transforms by
conjugation.

\item State how the projection formula changes if the generators have a
non-unit normalization.
\end{enumerate}
\end{exercise}
